\section{The Riccati equation and its two roots}
\label{bdc:sec:riccati}

Three probabilities sit close together in this process and are easily confused:
the cell may have avoided rupture, it may have died internally, or it may still
be productive. Their meanings differ, yet all three can be obtained from one
scalar differential equation. We derive that equation first, then follow its
flow from two initial conditions.

\subsection{First-step decomposition}
\label{bdc:sec:firststep}

Let
\[
I(t)=\pr{X_t\neq H \mid X_0=1}
\]
be the \emph{non-catastrophe probability}, the chance that rupture has not
occurred by time $t$. A path counted by $I(t)$ may still be productive or it may
already have died at $0$; it is excluded only if it has reached $H$.

Consider the first event experienced by the founding unit during a short
interval. Time homogeneity allows the process left after that event to be
followed for a further duration $t$:
\[
I(t+\Delta t) = \mu\,\Delta t
+ \bigl[1-(\lambda+\mu+\delta)\Delta t\bigr]I(t)
+ \lambda\,\Delta t\,I(t)^2 + o(\Delta t).
\]
If the founder dies, rupture has been avoided for good, giving the first term.
A rupture contributes nothing. If the founder divides, both descendant families
must avoid rupture, and the branching property gives probability $I(t)^2$.
After rearrangement,
\[
\frac{I(t+\Delta t)-I(t)}{\Delta t}
= \mu - (\lambda+\mu+\delta)I(t) + \lambda I(t)^2
+ \frac{o(\Delta t)}{\Delta t}.
\]
Taking $\Delta t\to0$ gives
\begin{equation}
\label{bdc:eq:riccati}
\ddt{I} = \mu - (\lambda+\mu+\delta)I + \lambda I^2,
\qquad I(0)=1.
\end{equation}
This is a Riccati equation, but its useful feature is simpler than the name
suggests: the right-hand side contains $I$ and the three rates, and nothing else.
The equation closes before any moments have been introduced.

\subsection{The two roots}
\label{bdc:sec:roots}

The quadratic on the right of \cref{bdc:eq:riccati} becomes transparent once its
roots are named.

\begin{lemma}[The characteristic quadratic has one root on either side of unity]
\label{bdc:lem:roots}
Let $\lambda>0$, $\mu\ge0$ and $\delta>0$, and set
\[
q(x)=\lambda x^2-(\lambda+\mu+\delta)x+\mu,
\qquad
\eta=\frac{\lambda+\mu+\delta}{2\lambda}.
\]
Then $q$ has two real roots
\[
a=\eta+\sqrt{\eta^2-\frac{\mu}{\lambda}},
\qquad
b=\eta-\sqrt{\eta^2-\frac{\mu}{\lambda}},
\]
satisfying $b<1<a$, with $ab=\mu/\lambda$ and
$a+b=(\lambda+\mu+\delta)/\lambda$. Moreover $b\in[0,1)$, with $b=0$ if and
only if $\mu=0$.
\end{lemma}

\begin{proof}
The roots are real and distinct because
\[
\eta^2-\frac{\mu}{\lambda}
=\frac{(\lambda+\mu+\delta)^2-4\lambda\mu}{4\lambda^2}
=\frac{(\lambda-\mu)^2+\delta(\delta+2\lambda+2\mu)}{4\lambda^2}>0
\]
when $\delta>0$. Since
$q(1)=\lambda-(\lambda+\mu+\delta)+\mu=-\delta<0$, the value at $1$ lies below
the axis of an upward-opening parabola, so one root lies on each side of unity.
Finally, Vieta's relation gives $ab=\mu/\lambda\ge0$. As $a>0$, this implies
$b\ge0$, with equality precisely when $\mu=0$.
\end{proof}

\begin{remark}[Standing assumptions]
\label{bdc:rem:assumptions}
Unless a limiting case is being considered, we take
$\lambda>0$, $\mu\ge0$ and $\delta>0$, and set $X_0=1$.
The case $\delta=0$ is the linear birth--death process of
\mextref{m:sec:bd}{Section 2.2.3 of \ChM}; it returns in
\cref{bdc:app:nocatastrophe} as a limit. The linear state dependence of the
rates also makes the process non-explosive, with finite moments at every finite
time. This permits the differentiations under expectation used in
\cref{bdc:sec:moments}.
\end{remark}

The Riccati equation now reads
\begin{equation}
\label{bdc:eq:riccatifact}
\ddt{I} = \lambda(I-a)(I-b).
\end{equation}
This factorisation gives a phase-line picture before it gives a formula.
The initial value $I(0)=1$ lies between $b$ and $a$, where the derivative is
negative. Consequently $I(t)$ decreases towards $b$ and never crosses it.
\Cref{bdc:fig:roots} draws this motion; the single observation
$q(1)=-\delta$ places the initial point between the two roots.

\begin{figure}[!htbp]
\centering
\includegraphics[width=0.72\textwidth]{F4_quadratic_roots.pdf}
\caption{The characteristic quadratic and its flow for $\lambda=1$, $\mu=0.2$
and $\delta=0.05$. The roots $b=0.188$ and $a=1.062$ lie on either side of
unity because $q(1)=-\delta$. A trajectory beginning at $I(0)=1$ moves
towards the stable root $b$. The distances $A=a-1$ and $B=1-b$ recur in the
closed forms.}
\label{bdc:fig:roots}
\end{figure}

Two short distances from unity make the later algebra easier to read:
\[
A=a-1,\qquad B=1-b.
\]
Vieta's relations then yield
\[
AB=(a-1)(1-b)=a+b-ab-1
=\frac{\lambda+\mu+\delta}{\lambda}-\frac{\mu}{\lambda}-1
=\frac{\delta}{\lambda},
\]
while $A+B=a-b$. We collect the identities here and name the exponential rate
that will recur:
\begin{equation}
\label{bdc:eq:identities}
ab=\frac{\mu}{\lambda},\qquad
a+b=\frac{\lambda+\mu+\delta}{\lambda},\qquad
AB=\frac{\delta}{\lambda},\qquad
A+B=a-b,\qquad
\theta=\lambda(a-b).
\end{equation}
The point of these relations is practical. Expressions which are cumbersome in
$\lambda$, $\mu$ and $\delta$ often fall into a recognisable shape in
$a$, $b$, $A$ and $B$.

\subsection{Solving for \texorpdfstring{$I(t)$}{I(t)}}
\label{bdc:sec:solveI}

Equation~\cref{bdc:eq:riccatifact} separates:
\[
\int \frac{\dd I}{(I-a)(I-b)} = \lambda t + \text{const}.
\]
Using
\[
\frac{1}{(I-a)(I-b)}
=\frac{1}{a-b}\lrb{\frac{1}{I-a}-\frac{1}{I-b}},
\]
we obtain
\begin{equation}
\label{bdc:eq:ratio}
\frac{I(t)-a}{I(t)-b}
=\frac{I(0)-a}{I(0)-b}\,\ee^{\lambda(a-b)t}
=\frac{I(0)-a}{I(0)-b}\,w,
\qquad w:=\ee^{\theta t}.
\end{equation}
For $I(0)=1$, the initial ratio is $-A/B$, and solving for $I(t)$ gives the
first main formula.

\begin{theorem}[The non-catastrophe probability]
\label{bdc:thm:I}
From one founder, the probability that catastrophe has not occurred by time $t$
is
\begin{equation}
\label{bdc:eq:I}
I(t) = \frac{aB + bA\,w}{B + A\,w},
\qquad w=\ee^{\theta t},\quad \theta=\lambda(a-b),
\end{equation}
where $a$ and $b$ are given in \cref{bdc:lem:roots}, with $A=a-1$ and
$B=1-b$.
\end{theorem}

As $t$ grows, $w$ grows and
\[
\lim_{t\to\infty}I(t)
=\lim_{t\to\infty}\frac{aB+bAw}{B+Aw}=b.
\]
When $\mu>0$, the paths which avoid rupture indefinitely are exactly those which
die internally first. Thus $b$ is the probability of internal extinction before
rupture. When $\mu=0$, death is impossible, $b=0$, and every path eventually
ruptures. The lower root is therefore both the stable point of the differential
equation and a probability. The upper root is not visited by $I(t)$, although it
will acquire its own probabilistic role through the burst-size law.

\subsection{The same equation again: internal extinction}
\label{bdc:sec:D}

Let
\[
D(t)=p_0(t)=\pr{X_t=0\mid X_0=1}
\]
be the probability of internal extinction by time $t$. The first-step argument
has the same shape as before. Death of the founder contributes certainty;
division requires both descendant families to die; rupture contributes nothing.
Consequently
\begin{equation}
\label{bdc:eq:riccatiD}
\ddt{D} = \mu - (\lambda+\mu+\delta)D + \lambda D^2,
\qquad D(0)=0.
\end{equation}
The meanings of $I$ and $D$ differ, but their motion is governed by the same
flow. For $b>0$, inserting $D(0)=0$ into \cref{bdc:eq:ratio} gives the initial
ratio $a/b$ and hence the following expression. When $b=0$, internal extinction
is impossible and the same formula is recovered directly as $D\equiv0$.

\begin{theorem}[The internal-extinction probability]
\label{bdc:thm:D}
The probability that the process has died out internally by time $t$ is
\begin{equation}
\label{bdc:eq:D}
D(t) = \frac{ab(w-1)}{aw-b}.
\end{equation}
\end{theorem}

The two solutions approach $b$ from opposite sides: $I$ descends from $1$ and
$D$ rises from $0$. Their narrowing separation is the probability that the cell
is still productive.

\subsection{The non-fixation probability}
\label{bdc:sec:Ifix}

Define
\[
\Ifix(t)=\pr{X_t\notin\{0,H\}\mid X_0=1}.
\]
Since $0$ and $H$ are the only ways to leave the productive states,
\[
\Ifix(t)=I(t)-D(t).
\]
Equivalently, if $T_0$ and $T_H$ are the two hitting times, then
$\Ifix(t)=\pr{t<T_0\wedge T_H}$. This is the survival event used for
conditioning in \fwd{singlecell}{Chapter~5}: the cell must remain intact
\emph{and} retain a positive load.

A direct first-step calculation is useful because it anticipates the
multi-founder case. Following a birth, the combined process is productive at
time $t$ provided no descendant family has triggered catastrophe and the two
families have not both died out. Those events have probabilities $I(t)^2$ and
$D(t)^2$, respectively, so the birth branch contributes $I^2-D^2$. Thus
\begin{equation}
\label{bdc:eq:ddIfix}
\ddt{\Ifix}
=-(\lambda+\mu+\delta)\Ifix+\lambda\lrb{I^2-D^2},
\qquad \Ifix(0)=1.
\end{equation}
Subtracting \cref{bdc:eq:riccatiD} from \cref{bdc:eq:riccati} gives the same
equation, as it should.

\begin{theorem}[The non-fixation probability]
\label{bdc:thm:Ifix}
The probability that the process remains productive at time $t$ is
\begin{equation}
\label{bdc:eq:Ifix}
\Ifix(t) = I(t)-D(t)
= \frac{aB+bAw}{B+Aw}-\frac{ab(w-1)}{aw-b}
= \frac{(a-b)^2\,w}{(B+Aw)(aw-b)}.
\end{equation}
\end{theorem}

The boundary values fit the interpretation. At $t=0$, the last expression is
one; as $t\to\infty$, its numerator grows as $w$ and its denominator as $w^2$,
so it tends to zero. Evaluating the differential equation at the origin gives
\[
\Ifix'(0)=-(\lambda+\mu+\delta)+\lambda=-(\mu+\delta),
\]
the combined rate at which the founding unit is lost either to death or to
rupture. A birth does not end productivity.

If $\mu=0$, then $b=0$ and $D\equiv0$, so $\Ifix=I$. When $\mu>0$, the
two functions separate, and conditioning on non-catastrophe alone admits cells
which are intact but already empty. \Cref{bdc:fig:fixation} shows this
separation and the eventual closing of the gap.

\begin{figure}[!htbp]
\centering
\includegraphics[width=0.72\textwidth]{F4_fate_probabilities.pdf}
\caption{Non-catastrophe, internal-extinction, and productive probabilities
for $\lambda=1$, $\mu=0.2$ and $\delta=0.05$. The non-catastrophe probability
$I(t)$ falls to $b$ while the internal-extinction probability $D(t)$ rises to
it. Their difference $\Ifix(t)$ is the probability of continued productivity
and tends to zero.}
\label{bdc:fig:fixation}
\end{figure}

\FloatBarrier

\subsection{Several founders}
\label{bdc:sec:multifounder}

Now suppose that the cell begins with $X_0=k$. Up to the first catastrophe, the
family trees and trigger clocks descending from the founders may be constructed
independently, as described in
\mextref{m:sec:bdc}{Section 2.2.4 of \ChM}. The chance that none of the $k$
families has triggered rupture by time $t$ is
\[
I_k(t)=\pr{X_t\neq H\mid X_0=k}=I(t)^k,
\]
whereas extinction of the whole population requires every family to die, giving
$D_k(t)=D(t)^k$. Their difference is
\begin{equation}
\label{bdc:eq:Ifixk}
\Ifix[,{k}](t)
=\pr{X_t\notin\{0,H\}\mid X_0=k}
=I(t)^k-D(t)^k.
\end{equation}

\begin{remark}[Non-fixation does not factorise]
\label{bdc:rem:multifounder}
For $\mu>0$, the tempting identity
$\Ifix[,{k}]=\Ifix^k$ is false. If one family dies while another remains
productive, the cell has not fixed; these mixed configurations are retained by
$I^k-D^k$ but omitted by $(I-D)^k$. Catastrophe is different from the death of
one lineage: the first trigger ruptures the cell and stops all lineages at once.
Only when $\mu=0$, so that $D\equiv0$, do the two expressions coincide.
\end{remark}

\Cref{bdc:fig:multifounder} makes the omitted configurations visible. For two
founders the naive power lies below the exact probability, and the gap is the
probability that one lineage has died while the other is still productive.

\begin{figure}[!htbp]
\centering
\includegraphics[width=0.94\textwidth]{F4_multi_founder.pdf}
\caption{Non-fixation from several founders. (a) Exact probabilities
$\Ifix[,{k}](t)=I(t)^k-D(t)^k$ for $k=1,2,3,5$. (b) For two founders,
$[I(t)-D(t)]^2$ lies below the exact result because it omits mixed configurations
with one extinct and one productive lineage. $\lambda=1$, $\mu=0.2$ and
$\delta=0.05$.}
\label{bdc:fig:multifounder}
\end{figure}

These multi-founder quantities return in \fwd{singlecell}{Chapter~5}, where
the generating functions provide a second derivation and allow several incoming
particles to found one infection.

\FloatBarrier
